% ===== §5 The Third Cause: The Arrest of Geometric-Phase Accumulation =====
\section{The Third Cause: The Arrest of Phase}
\label{p17:sec:phase}

\noindent
The first cause lay in the code, the second in the act of prediction. The third lies in neither: it lies in the \emph{value} that the coupling was supposed to generate, and in the geometry of its generation. Two subjects may share a grammar and attribute their errors faithfully and still find that nothing new is any longer being made between them---that the bond, though intact, has stopped \emph{generating}. To say precisely what has failed, and why ``passing without meeting'' is the exact phrase for it, the section draws on the geometric apparatus the series developed to distinguish good cycles from bad. This is the layer at which the descriptive image of \xref{p17:sub:symptom-three}---estrangement as \emph{zero-phase adjacency}---receives its formal content.

\subsection{Value as holonomy: a recapitulation}
\label{p17:sub:phase-recap}

The series has argued that the value generated by a relational coupling is not a substance stored in either party but a \emph{geometric phase}: a quantity that accrues only along a closed path traversed \emph{together}, and that depends on the path and not merely on its endpoints.\footnote{For the holonomy apparatus---value as the geometric phase accumulated around a closed interpretive loop, the curvature of the relational connection as the local density of value-generation, and the sign of the holonomy as the formal criterion separating a good (positive) cycle from an extractive (zero or negative) one---see \textcite{paperix}. The present section uses this apparatus to characterise the failure of accumulation.} Two subjects, interpreting each other and the world, trace a loop in the space of shared meaning; if the loop closes upon a configuration that is not identical to its start but \emph{lifted}---returned with a surplus that neither held before and neither can claim alone---then the cycle has positive holonomy, and that surplus is the value generated. A bad cycle is one that returns to exactly where it began, having moved meaning nowhere (zero holonomy), or one that returns depleted, having drained one party to inflate the other (negative holonomy). On this account value is intrinsically \emph{path-dependent and joint}: it cannot be accumulated by one subject alone, and it cannot be accumulated at all unless a loop is actually closed.

\subsection{The geometry of passing-without-meeting}
\label{p17:sub:phase-geometry}

This apparatus yields an exact definition of the estrangement the paper has so far only described. \emph{Surechigai} is the condition in which two trajectories are brought close---adjacent in the space of meaning, near enough to touch---while the phase that would accumulate \emph{between} them goes to zero. The two are proximate and the holonomy is nil. This is why estrangement is so precisely \emph{not} collision: a collision is an interaction, an exchange of momentum, an event with consequence; zero-phase adjacency is the absence of consequence at the very point of nearness. The subjects are close enough that something should be generated, and nothing is. They pass.

\begin{claim}[The arrest of phase]
\emph{Surechigai}, at this layer, is the \emph{arrest of geometric-phase accumulation}: the interpretive loop that the two subjects were tracing together bifurcates or fails to close, so that no surplus is lifted and no value is generated, even as the two remain adjacent. Estrangement is not the divergence of two paths into the distance---that would be mere separation---but their nearness \emph{without} accumulation: zero holonomy at close approach. The bond is intact; the generation has stopped.
\end{claim}

\subsection{Three forms of arrest}
\label{p17:sub:phase-forms}

The arrest takes three forms, which are the three causes this layer contributes to the etiology, and which correspond to three ways a jointly traced loop can fail to generate.

The first is \emph{bifurcation of the loop}. The two subjects, formerly tracing one loop in shared meaning, begin to trace two---interpreting the same events along paths that no longer coincide, so that there is no single closed loop around which a common phase could accumulate. Each may be generating value on their own path; but value, on the holonomy account, is what accrues \emph{jointly}, around a shared loop, and two private loops, however active, generate no \emph{relational} surplus. This is the value-theoretic form of the predictive withdrawal of \xref{p17:sec:prediction}: where prediction has disinvested, the loop has split.

The second is \emph{the loop that fails to close}. Here the two are still on one path, but the path no longer returns---each interpretive venture is left open, unresolved, never brought back to be witnessed and lifted into shared surplus. The series has argued that the good cycle is one that \emph{refuses settlement} while nonetheless closing its loop, so that meaning is moved without being fixed; the failure here is the opposite-yet-adjacent pathology of a loop that, never closing at all, accumulates nothing because accumulation requires return. Activity without closure is not generativity; it is dispersion.

The third is \emph{accumulated negative curvature}. The path may close, but it now runs across a region of the relational manifold that earlier passings have bent. Past estrangements are not erased; they deposit curvature---scar tissue in the geometry---and a loop traced across scarred ground accumulates phase less readily, or accumulates it with the wrong sign. This is the layer's contribution to the paper's account of \emph{time} (\xref{p17:sec:dissipation}, \xref{p17:sec:histmat}): a bond keeps a ledger, and the geometry through which today's loop must pass is the record of every loop traced before. Estrangement, here, is partly the debt of estrangements past.

\subsection{Formalising the ``objective obstacle''}
\label{p17:sub:phase-obstacle}

This layer also lets the paper make good on a phrase used loosely until now: the \emph{objective obstacle} that intervenes in the emergence of value. The obstacle is objective in the precise sense that it is not a matter of either subject's will or feeling but a feature of the \emph{geometry} they jointly inhabit. A loop can fail to close, bifurcate, or run across negative curvature regardless of how much either party wishes it otherwise, just as a path on a curved surface accumulates a holonomy fixed by the surface and not by the traveller's intentions. This is what it means to say that the arrest of value is not, at bottom, a psychological failure but a geometric condition---and it is what makes the arrest, like the two causes before it, in the first instance \emph{blameless}. No one need have done wrong for a loop to fail to close; the curvature was there, or the paths bifurcated, and the phase did not accumulate.

\begin{claim}[The objective obstacle]
The arrest of value-generation is an \emph{objective} obstacle: a property of the relational geometry the two subjects jointly inhabit, not of either subject's will or feeling. A loop fails to close, splits, or crosses negative curvature as a feature of the manifold, not as an act of either party. Hence the intervention of obstacle in the emergence of value is, like encoding mismatch and unattributable error, prior to responsibility---and hence, too, it cannot be removed by good will alone, but only by the slow re-shaping of the geometry through which the loops are traced.
\end{claim}

\subsection{Closing the internal causes}
\label{p17:sub:phase-close}

With the arrest of phase the etiology has traced estrangement through the three layers internal to the dyad: the code (\xref{p17:sec:encoding}), the prediction (\xref{p17:sec:prediction}), and the value (\xref{p17:sec:phase}). They are not three separate troubles but three depths of one: a mismatch of grammars makes errors harder to attribute; unattributable error triggers predictive withdrawal; predictive withdrawal splits the loop and arrests the phase. Each layer feeds the next, and each, taken alone, is blameless---a property of the pairing, the channel, or the geometry, and not yet of any will. This blamelessness is the great finding of the internal etiology, and it is what the ethics of \xref{p17:sec:asymmetry} will have to reckon with: that the commonest and deepest forms of estrangement arrive before responsibility does, and that the reach for blame is itself, so often, the first genuine wrong.

But the dyad is not closed, and the three internal causes do not exhaust the field. Around the two subjects run the powers the introduction named---families, scripts, institutions, the sediment of histories---and these do not merely add a fourth cause alongside the three. They change the standing of the first three, deciding which grammar is default, injecting the noise that makes error unattributable, bending the manifold across which the loops must run. To take their measure, the paper must now widen from the dyad to the field, and for that it requires a theory of how opaque subjects build, or fail to build, order between them under no common authority. It requires, that is, a reconstruction of the theory of international relations---to which the next section turns.
